\section{C++/Tcl Glue Code Template Framework}
\label{app:cpp}

Complete source for the three-header template framework described in
Section~\ref{sec:impl}. Requires C++17 and Tcl~8.6 C~API headers.

\subsection*{D.1 \texttt{tcl/type.h} --- Custom \texttt{Tcl\_ObjType} Registration}

\begin{lstlisting}[language=C++, caption={\texttt{tcl/type.h} --- type registration template}]
#pragma once
#include <tcl.h>
#include <exception>
#include <string>
#include <sstream>
#include <cstring>

namespace tcl {

template <typename T>
class Type {
public:
    static Tcl_ObjType* GetType() {
        static std::string nameStorage(TypeName());
        static Tcl_ObjType type = {
            nameStorage.c_str(),
            FreeInternalRep,
            DupInternalRep,
            UpdateString,
            SetFromAny
        };
        static bool registered = false;
        if (!registered) { Tcl_RegisterObjType(&type); registered = true; }
        return &type;
    }

    static Tcl_Obj* New(T* value) {
        Tcl_Obj* obj = Tcl_NewObj();
        if (!obj) throw std::runtime_error("Failed to allocate Tcl object");
        obj->bytes = NULL;
        obj->typePtr = GetType();
        obj->internalRep.otherValuePtr = static_cast<void*>(value);
        return obj;
    }

    static void FreeInternalRep(Tcl_Obj* obj) {
        delete static_cast<T*>(obj->internalRep.otherValuePtr);
        obj->internalRep.otherValuePtr = nullptr;
    }

    static void DupInternalRep(Tcl_Obj* src, Tcl_Obj* dup) {
        dup->internalRep.otherValuePtr =
            new T(*static_cast<T*>(src->internalRep.otherValuePtr));
        dup->typePtr = GetType();
    }

    static void UpdateString(Tcl_Obj* obj) {
        T* value = static_cast<T*>(obj->internalRep.otherValuePtr);
        std::string str = ToString(*value);
        obj->bytes = Tcl_Alloc(str.size() + 1);
        strcpy(obj->bytes, str.c_str());
        obj->length = str.size();
    }

    static int SetFromAny(Tcl_Interp* interp, Tcl_Obj* obj) {
        try {
            T value = FromAny(interp, obj);
            if (obj->typePtr && obj->typePtr->freeIntRepProc)
                obj->typePtr->freeIntRepProc(obj);
            obj->internalRep.otherValuePtr = new T(value);
            obj->typePtr = GetType();
            Tcl_InvalidateStringRep(obj);
            return TCL_OK;
        } catch (const std::exception& e) {
            if (interp) Tcl_SetObjResult(interp, Tcl_NewStringObj(e.what(), -1));
            return TCL_ERROR;
        }
    }

    static T* GetInternalRep(Tcl_Interp* interp, Tcl_Obj* obj) {
        if (obj->typePtr != GetType()) {
            if (SetFromAny(interp, obj) != TCL_OK) {
                std::ostringstream oss;
                oss << "Failed to convert to \"" << TypeName() << "\"";
                throw std::runtime_error(oss.str());
            }
        }
        return static_cast<T*>(obj->internalRep.otherValuePtr);
    }

    // Default ToString --- throws; specialize to enable `puts $obj`
    static std::string ToString(const T&) {
        std::ostringstream oss;
        oss << "Type \"" << TypeName() << "\" can't be cast to string";
        throw std::runtime_error(oss.str());
    }

    // Default FromAny falls back to string conversion
    static T FromAny(Tcl_Interp*, Tcl_Obj* const obj) {
        return FromString(Tcl_GetString(obj));
    }

    static T FromString(const std::string&) {
        std::ostringstream oss;
        oss << "Type \"" << TypeName() << "\" can't be cast from string";
        throw std::runtime_error(oss.str());
    }

private:
    static const char* TypeName();  // Must be specialized per type
};

} // namespace tcl
\end{lstlisting}

\begin{table}[H]
\centering
\caption{Specialization points for \texttt{tcl::Type<T>}}
\begin{tabular}{lll}
\toprule
\textbf{Function} & \textbf{Required} & \textbf{Purpose} \\
\midrule
\texttt{TypeName()} & Yes & Unique name for Tcl's type registry \\
\texttt{ToString(const T\&)} & No & Serializes \texttt{T} for \texttt{puts \$obj} \\
\texttt{FromAny(interp, obj)} & No & Parses a \texttt{Tcl\_Obj} into \texttt{T} \\
\texttt{FromString(const string\&)} & No & Fallback used by default \texttt{FromAny} \\
\bottomrule
\end{tabular}
\end{table}

\subsection*{D.2 \texttt{tcl/obj\_cast.h} --- Bidirectional Type Casting}

\begin{lstlisting}[language=C++, caption={\texttt{tcl/obj\_cast.h} (excerpt showing key specializations)}]
#pragma once
#include <tcl.h>
#include <string>
#include <stdexcept>
#include <type_traits>
#include "tcl/type.h"

namespace tcl { namespace obj_cast {

    template <typename T>
    Tcl_Obj* from(const T&) {
        static_assert(!std::is_same_v<T, T>,
            "No tcl::obj_cast::from<T> specialization for this type.");
        return nullptr;
    }

    template <typename T>
    T to(Tcl_Interp* interp, Tcl_Obj* const obj) {
        if constexpr (std::is_pointer_v<T>)
            return tcl::Type<std::remove_pointer_t<T>>::GetInternalRep(interp, obj);
        else
            return *tcl::Type<T>::GetInternalRep(interp, obj);
    }

    // int
    template <> inline Tcl_Obj* from<int>(const int& v) { return Tcl_NewIntObj(v); }
    template <> inline int to<int>(Tcl_Interp* i, Tcl_Obj* const o) {
        int v; Tcl_GetIntFromObj(i, o, &v); return v; }

    // double
    template <> inline Tcl_Obj* from<double>(const double& v) { return Tcl_NewDoubleObj(v); }
    template <> inline double to<double>(Tcl_Interp* i, Tcl_Obj* const o) {
        double v; Tcl_GetDoubleFromObj(i, o, &v); return v; }

    // std::string
    template <> inline Tcl_Obj* from<std::string>(const std::string& v) {
        return Tcl_NewStringObj(v.c_str(), -1); }
    template <> inline std::string to<std::string>(Tcl_Interp*, Tcl_Obj* const o) {
        int len; const char* s = Tcl_GetStringFromObj(o, &len);
        return std::string(s, len); }

    // bool
    template <> inline Tcl_Obj* from<bool>(const bool& v) { return Tcl_NewBooleanObj(v); }
    template <> inline bool to<bool>(Tcl_Interp* i, Tcl_Obj* const o) {
        int v; Tcl_GetBooleanFromObj(i, o, &v); return static_cast<bool>(v); }

    // Tcl_Obj* passthrough
    template <> inline Tcl_Obj* from<Tcl_Obj*>(Tcl_Obj* const& v) { return v; }
    template <> inline Tcl_Obj* to<Tcl_Obj*>(Tcl_Interp*, Tcl_Obj* const o) { return o; }

}} // namespace tcl::obj_cast
\end{lstlisting}

\subsection*{D.3 \texttt{tcl/cpp\_api.h} --- Glue Macros and Template Functions}

\begin{lstlisting}[language=C++, caption={Key macros from \texttt{tcl/cpp\_api.h}}]
// Constructor: new T(arg1, arg2, ...)
#define TCLCPPG_CREATE_CMD(interp, cmdName, className, ...) \
    Tcl_CreateObjCommand(interp, cmdName, \
        [](ClientData cd, Tcl_Interp *ti, int oc, Tcl_Obj *const ov[]) -> int { \
            return tcl::cpp_api::create<className, __VA_ARGS__>(cd, ti, oc, ov); \
        }, NULL, NULL)

// Getter (no extra args): object by value
#define TCLCPPG_GETTER_CMD_NOARGS(interp, cmdName, className, returnType, methodName) \
    Tcl_CreateObjCommand(interp, cmdName, \
        [](ClientData cd, Tcl_Interp *ti, int oc, Tcl_Obj *const ov[]) -> int { \
            return tcl::cpp_api::getter<className, returnType>( \
                cd, ti, oc, ov, &className::methodName); \
        }, NULL, NULL)

// Setter (with extra args): object by variable name
#define TCLCPPG_SETTER_CMD(interp, cmdName, className, returnType, methodName, ...) \
    Tcl_CreateObjCommand(interp, cmdName, \
        [](ClientData cd, Tcl_Interp *ti, int oc, Tcl_Obj *const ov[]) -> int { \
            return tcl::cpp_api::setter<className, returnType, __VA_ARGS__>( \
                cd, ti, oc, ov, &className::methodName); \
        }, NULL, NULL)
\end{lstlisting}

\begin{table}[H]
\centering
\caption{Summary of \texttt{TCLCPPG\_*} macros}
\begin{tabular}{@{}L{5.8cm}L{3.0cm}L{5.2cm}@{}}
\toprule
\textbf{Macro} & \textbf{Tcl convention} & \textbf{C++ side} \\
\midrule
\texttt{TCLCPPG\_CREATE\_CMD} & \texttt{name arg1 arg2 ...} & \texttt{new Class(args)} \\
\texttt{TCLCPPG\_GETTER\_CMD\_NOARGS} & \texttt{name object} & \texttt{const} member, no extra args \\
\texttt{TCLCPPG\_GETTER\_CMD} & \texttt{name object arg1 ...} & \texttt{const} member with extra args \\
\texttt{TCLCPPG\_SETTER\_CMD\_NOARGS} & \texttt{name objectVar} & non-\texttt{const} member, no extra args \\
\texttt{TCLCPPG\_SETTER\_CMD} & \texttt{name objectVar arg1 ...} & non-\texttt{const} member with extra args \\
\texttt{TCLCPPG\_FREE\_CMD} & \texttt{name arg1 arg2 ...} & free function or static method \\
\bottomrule
\end{tabular}
\end{table}

\subsection*{D.4 Complete Example: \texttt{CppVooPoint} (\texttt{point.cpp})}

\begin{lstlisting}[language=C++, caption={Step 1--2: C++ class and \texttt{tcl::Type} specialization}]
#include <tcl.h>
#include <string>
#include <cmath>
#include <sstream>
#include "tcl/cpp_api.h"

class VooPoint {
    double m_x, m_y;
    std::string m_name;
    int m_id;
    bool m_active;
public:
    VooPoint(double x, double y, const std::string& name, int id, bool active)
        : m_x(x), m_y(y), m_name(name), m_id(id), m_active(active) {}
    VooPoint() : m_x(0.0), m_y(0.0), m_name("point"), m_id(0), m_active(true) {}

    double             getX()      const { return m_x; }
    double             getY()      const { return m_y; }
    const std::string& getName()   const { return m_name; }
    int                getId()     const { return m_id; }
    bool               getActive() const { return m_active; }
    void setX(double v)            { m_x = v; }
    void setY(double v)            { m_y = v; }
    void setName(std::string v)    { m_name = std::move(v); }
    void setId(int v)              { m_id = v; }
    void setActive(bool v)         { m_active = v; }
    double distance() const { return std::sqrt(m_x*m_x + m_y*m_y); }
};

namespace tcl {
template <> const char* Type<VooPoint>::TypeName() { return "VooPoint"; }
template <> std::string Type<VooPoint>::ToString(const VooPoint& p) {
    std::ostringstream oss;
    oss << p.getX() << " " << p.getY() << " "
        << p.getName() << " " << p.getId() << " " << p.getActive();
    return oss.str();
}
template <> VooPoint Type<VooPoint>::FromAny(Tcl_Interp* interp, Tcl_Obj* const obj) {
    int objc; Tcl_Obj** objv;
    if (Tcl_ListObjGetElements(interp, obj, &objc, &objv) != TCL_OK || objc != 5)
        throw std::runtime_error("Expected list of 5 elements: x y name id active");
    return VooPoint(
        obj_cast::to<double>(interp, objv[0]),
        obj_cast::to<double>(interp, objv[1]),
        obj_cast::to<std::string>(interp, objv[2]),
        obj_cast::to<int>(interp, objv[3]),
        obj_cast::to<bool>(interp, objv[4]));
}
}
\end{lstlisting}

\begin{lstlisting}[language=C++, caption={Step 3--4: default constructor helper and package initializer}]
// Zero-argument default constructor
static int VooPoint_newDefault(ClientData, Tcl_Interp* interp,
                               int objc, Tcl_Obj* const objv[]) {
    if (objc != 1) { Tcl_WrongNumArgs(interp, 1, objv, ""); return TCL_ERROR; }
    try {
        Tcl_SetObjResult(interp, tcl::Type<VooPoint>::New(new VooPoint()));
    } catch (const std::exception& e) {
        Tcl_SetObjResult(interp, Tcl_NewStringObj(e.what(), -1));
        return TCL_ERROR;
    }
    return TCL_OK;
}

extern "C" { int Point_Init(Tcl_Interp* interp); }

int Point_Init(Tcl_Interp* interp) {
    Tcl_CreateNamespace(interp, "::CppVooPoint", NULL, NULL);

    // Constructors
    TCLCPPG_CREATE_CMD(interp, "::CppVooPoint::new",
                       VooPoint, double, double, std::string, int, bool);
    Tcl_CreateObjCommand(interp, "::CppVooPoint::new()",
                         VooPoint_newDefault, NULL, NULL);

    // Getters
    TCLCPPG_GETTER_CMD_NOARGS(interp, "::CppVooPoint::get.x",     VooPoint, double,             getX);
    TCLCPPG_GETTER_CMD_NOARGS(interp, "::CppVooPoint::get.y",     VooPoint, double,             getY);
    TCLCPPG_GETTER_CMD_NOARGS(interp, "::CppVooPoint::get.name",  VooPoint, const std::string&, getName);
    TCLCPPG_GETTER_CMD_NOARGS(interp, "::CppVooPoint::get.id",    VooPoint, int,                getId);
    TCLCPPG_GETTER_CMD_NOARGS(interp, "::CppVooPoint::get.active",VooPoint, bool,               getActive);
    TCLCPPG_GETTER_CMD_NOARGS(interp, "::CppVooPoint::distance",  VooPoint, double,             distance);

    // Setters
    TCLCPPG_SETTER_CMD(interp, "::CppVooPoint::set.x",     VooPoint, void, setX,     double);
    TCLCPPG_SETTER_CMD(interp, "::CppVooPoint::set.y",     VooPoint, void, setY,     double);
    TCLCPPG_SETTER_CMD(interp, "::CppVooPoint::set.name",  VooPoint, void, setName,  std::string);
    TCLCPPG_SETTER_CMD(interp, "::CppVooPoint::set.id",    VooPoint, void, setId,    int);
    TCLCPPG_SETTER_CMD(interp, "::CppVooPoint::set.active",VooPoint, void, setActive,bool);

    return Tcl_PkgProvideEx(interp, "VooPointCpp", "1.0.0", NULL);
}
\end{lstlisting}

\begin{table}[H]
\centering
\caption{Summary of migration steps for VOO C++}
\begin{tabular}{L{4.0cm}L{2.5cm}L{6.5cm}}
\toprule
\textbf{Step} & \textbf{Location} & \textbf{What it does} \\
\midrule
C++ class  & \texttt{point.cpp} & Plain C++ --- no Tcl headers required \\
\texttt{TypeName()} & \texttt{tcl::Type<>} specialization & Registers \texttt{VooPoint} in Tcl's type registry \\
\texttt{ToString()} & \texttt{tcl::Type<>} specialization & Enables \texttt{puts \$obj} and list/dict storage \\
\texttt{FromAny()}  & \texttt{tcl::Type<>} specialization & Parses a 5-element Tcl list back into a \texttt{VooPoint} \\
\texttt{VooPoint\_newDefault} & standalone command & Zero-argument constructor (\texttt{new()}) \\
\texttt{Point\_Init} & package init entry & Registers all Tcl commands; called by \texttt{load} \\
\bottomrule
\end{tabular}
\end{table}

The resulting Tcl API is identical to the pure-Tcl VOO declaration --- only the namespace
prefix (\texttt{CppVooPoint::} vs.\ \texttt{VooPoint::}) differs, so no caller code
needs to change when migrating from Tcl to C++.
